\begin{frame}{자주 하는 실수: 카운트 방식 선택}
  코드의 \texttt{for w in set(email)} 는 같은 단어가 여러 번 나와도
  \textbf{한 번만} 센다(베르누이) --- 실전에선 이 선택 자체가 버그의
  원인:
  \begin{itemize}
    \item `free' 10번 메일은 분명 1번 메일보다 스팸 가능성이 높은데,
      \texttt{set(email)} 은 \textbf{동일하게 처리} --- 신호 9개를
      버림.
    \item 반대로, \textbf{단어 수로 정규화하지 않은 채} 횟수만 곱해
      버리면(다항 + 스무딩 없음), \textbf{긴 메일일수록} 모든
      확률 곱이 무조건 작아져서 \textbf{문서 길이가 예측에 영향}
      을 주는 엉뚱한 인과가 생김.
    \item ``등장 \textbf{여부만} 볼 것인가(베르누이), \textbf{횟수}
      까지 볼 것인가(다항), 볼 때 어떻게 정규화할 것인가''를
      \textbf{의식적으로} 정하고, 데이터셋 특성(길이가 신호인가,
      반복 빈도가 신호인가)에 맞게 고를 것.
  \end{itemize}
\end{frame}
